\chapter{Experimental Results}

\setlength{\parindent}{1em}

\begin{quote}
\small\itshape
Note on what is measured versus planned. Every table in this chapter reports
real measurements from the deployed system: Whisper \texttt{small} paired with
Qwen2.5-7B-Instruct, run on a CPU-only reference machine with
\texttt{temperature}\,$=0$. These figures were obtained by running
\texttt{backend/evaluation/run\_stt\_eval.py} and \texttt{run\_extraction\_eval.py}
against the full 274-row corpus and a 6-clip recorded-audio pilot. The three
\textbf{ablation studies} (Whisper size, LLM size, few-shot vs.\ zero-shot,
Section~5.5) compare against alternative configurations that were not run in
this pass---each needs a container restart with a different model and a
30-minute re-run of the full corpus---so that section states the exact
comparison design and the one data point already measured (the deployed
configuration) for each, rather than presenting incomplete tables.
\end{quote}

\section{Baseline / Evaluation Setup}

\subsection{Evaluation Corpus and Conditions}

All experiments use the labeled corpus of Section~4.1: 274 utterances
(111 Arabic, 110 English, 53 code-switched), comprising 294 gold transactions,
of which 19 utterances are multi-transaction. Two evaluation \emph{conditions}
isolate different parts of the pipeline:

\begin{description}[leftmargin=1.4em, style=nextline]
    \item[Condition (a) --- Gold transcript $\rightarrow$ extraction.] The gold
    transcript is sent directly to the text-extraction endpoint, bypassing speech
    recognition. This isolates the language model and post-processing from any
    transcription noise and measures the \emph{ceiling} of extraction quality.
    Run over all 274 utterances.
    \item[Condition (b) --- Audio $\rightarrow$ STT $\rightarrow$ extraction.] A
    recorded audio clip is sent through the full voice endpoint. This measures
    \emph{end-to-end} quality, including the compounding of transcription errors
    into extraction errors. Run over the 6-clip recorded-audio pilot (four
    Arabic, two English) described in Section~4.1; this is a small qualitative
    pilot, not a statistically powered benchmark, and is reported as such
    throughout this chapter.
\end{description}

\subsection{Reference Hardware}

The reference deployment is the CPU-only configuration of NFR-06: a single
machine with 16\,GB RAM and no GPU, running the three Docker Compose services.
The speech model is Whisper \texttt{small} (\texttt{int8}) and the language
model is Qwen2.5-7B-Instruct served by Ollama with \texttt{temperature}\,$=0$
and schema-constrained decoding. All latency figures are wall-clock times
measured on this reference machine. Determinism was confirmed directly: an
independent full re-run of condition (a) reproduced identical aggregate
accuracy figures to four significant digits, as expected from
temperature-0 decoding.

\paragraph{Note on system version.} The corpus and results reported in this
chapter were measured against the seventeen-category, Saudi-Arabia-first
(SAR-default) configuration of the extraction schema described earlier in
this thesis. The system has since been extended for the Egyptian market
(Section~4.5.1): the category taxonomy grew from seventeen to twenty-nine
keys and the default currency changed from SAR to EGP. The evaluation
corpus itself is already Egyptian-dialect-heavy and EGP-denominated (Table~4.1),
so the measured accuracy figures and the two root causes identified in
Section~5.7 remain representative of the system's behavior; only the
category-taxonomy size referenced in Section~5.7.3 reflects the
seventeen-key schema in place at evaluation time. Re-running the full
evaluation against the twenty-nine-key schema is listed in Chapter~6 as an
immediate next step, alongside the ablation studies. The corpus has also grown
since this evaluation was run: eight further Arabic utterances exercising the
new Egyptian categories directly (\texttt{internet\_mobile}, \texttt{salary},
\texttt{coffee}, \texttt{subscriptions}, \texttt{savings}, \texttt{clothes},
\texttt{transfer\_out}, \texttt{freelance}) were added afterward and are
included in the corpus figures of Table~4.1 (282 rows) but are \emph{not} part
of the 274-utterance, 294-transaction sample the numbers in this chapter
describe; they are folded into the same future re-evaluation pass.

\section{Metrics}

\subsection{Word Error Rate (Speech-to-Text)}

Transcription quality is measured with Word Error Rate (WER), the standard ASR
metric~\cite{morris2004wer}, computed with the \texttt{jiwer}
library~\cite{jiwer}:
\[
\text{WER} = \frac{S + D + I}{N},
\]
where $S$, $D$, and $I$ are the number of substituted, deleted, and inserted
words required to transform the hypothesis into the reference, and $N$ is the
number of words in the reference. Because Arabic orthography varies, Arabic
references and hypotheses are normalized before scoring---diacritics are stripped,
alef variants are unified, and the \emph{ta-marbuta}/\emph{ha} ending is
normalized---following the project's \texttt{normalize\_ar.py}; Character Error
Rate (CER) is additionally reported for Arabic since WER can be a harsh metric
on short utterances with heavy dialectal variation.

\subsection{Extraction Accuracy and F\textsubscript{1}}

Predicted transactions are first aligned to gold transactions by greedy
best-match on \texttt{(transaction\_type, amount)}; unmatched predictions count
as false positives and unmatched gold transactions as false negatives. Over the
matched pairs:

\begin{itemize}[leftmargin=1.4em]
    \item \texttt{amount}, \texttt{transaction\_type}, \texttt{currency}, and
    \texttt{date} are scored by \textbf{exact-match accuracy} (amount within a
    1\% tolerance);
    \item \texttt{category} and \texttt{name} are scored by
    \textbf{precision, recall, and F\textsubscript{1}}~\cite{powers2011fmeasure}
    (name matched case-insensitively after normalization), with
    \[
    P=\frac{TP}{TP+FP}, \quad R=\frac{TP}{TP+FN}, \quad F_1=\frac{2PR}{P+R};
    \]
    \item a \textbf{transaction-count accuracy} reports how often the model
    detected the correct \emph{number} of transactions in an utterance.
\end{itemize}

\section{Speech-to-Text Results}

Table~\ref{tab:wer} reports WER on the 6-clip recorded-audio pilot. The Arabic
target is WER\,$\leq$\,20\% and the English target is WER\,$\leq$\,12\%
(NFR-02). \textbf{With only 4 Arabic and 2 English clips, these figures are a
qualitative signal, not a powered estimate}---each is a single wrong/right word
away from moving several percentage points--- but the direction is clear and
consistent with the extraction-stage findings below.

\begin{table}[H]
\centering
\caption{Word Error Rate on the 6-clip recorded-audio pilot (Whisper \texttt{small}, \texttt{int8}). Lower is better.}
\label{tab:wer}
\small
\renewcommand{\arraystretch}{1.3}
\begin{tabularx}{\textwidth}{X c c c c}
\toprule
\textbf{Language subset} & \textbf{WER (\%)} & \textbf{CER (\%)} & \textbf{Target WER (\%)} & \textbf{Clips} \\
\midrule
Arabic (\texttt{ar})  & 76.0 & 44.7 & $\leq 20$ & 4 \\
English (\texttt{en}) & 17.4 & ---  & $\leq 12$ & 2 \\
\bottomrule
\end{tabularx}
\end{table}

The English WER (17.4\%, $n=2$) is close to, but above, the 12\% target; the
Arabic WER (76.0\%, $n=4$) is far above the 20\% target. Two of the four Arabic
clips were the team's own re-takes of a single sentence in Egyptian colloquial
Arabic recorded under informal (non-studio) conditions, and Whisper's transcripts
of them were badly garbled (in one case the language was even mis-detected as
Japanese). This is consistent with the known difficulty dialectal, informally
recorded Arabic poses for general-purpose ASR (Section~2.1) and is discussed
further in Section~5.7.

\section{Extraction Results}

\subsection{Per-Field Accuracy}

Table~\ref{tab:extract-main} reports per-field extraction quality under both
conditions. The targets from NFR-02 are amount $\geq$\,90\%, type $\geq$\,85\%,
and category $\geq$\,80\%.

\begin{table}[H]
\centering
\caption{Per-field extraction quality. Exact-match accuracy for amount, type, currency, date; F\textsubscript{1} for category and name.}
\label{tab:extract-main}
\small
\renewcommand{\arraystretch}{1.3}
\begin{tabularx}{\textwidth}{X c c c}
\toprule
\textbf{Field} & \textbf{(a) Gold ($n$=274)} & \textbf{(b) End-to-end ($n$=6)} & \textbf{Target} \\
\midrule
\texttt{amount} (acc.)            & 65.6\% & 100.0\% & $\geq 90\%$ \\
\texttt{transaction\_type} (acc.) & 99.3\% & 100.0\% & $\geq 85\%$ \\
\texttt{currency} (acc.)          & 96.6\% & 71.4\%  & --- \\
\texttt{date} (acc.)              & 99.3\% & 100.0\% & --- \\
\texttt{category} (F\textsubscript{1}) & 0.554 & 0.222 & $\geq 0.80$ \\
\texttt{name} (F\textsubscript{1})     & 0.675 & 0.333 & --- \\
\midrule
Transaction-count acc.            & 100.0\% & 100.0\% & --- \\
\bottomrule
\end{tabularx}
\end{table}

\texttt{transaction\_type}, \texttt{date}, and transaction-count accuracy all
meet or substantially exceed their targets. \texttt{amount} and
\texttt{category} fall well short of theirs in the gold condition---65.6\%
against a 90\% target for amount, and an F\textsubscript{1} of 0.554 against an
0.80 target for category. Section~5.7 shows these are not diffuse, unexplained
weaknesses: both trace to specific, identifiable causes with a clear fix.

\subsection{Per-Language Breakdown}

Table~\ref{tab:extract-lang} breaks extraction accuracy down by language subset
(condition (a)). The gap is stark for \texttt{amount} and is the single most
important result in this chapter for a bilingual system.

\begin{table}[H]
\centering
\caption{Extraction accuracy by language subset (condition (a), gold transcript, $n$=111/110/53).}
\label{tab:extract-lang}
\small
\renewcommand{\arraystretch}{1.3}
\begin{tabularx}{\textwidth}{X c c c c c}
\toprule
\textbf{Subset} & \textbf{type} & \textbf{amount} & \textbf{currency} & \textbf{date} & \textbf{category F\textsubscript{1}} \\
\midrule
Arabic        & 99.2\%  & \textbf{37.3\%} & 98.3\% & 98.3\%  & 0.564 \\
English       & 99.2\%  & \textbf{99.2\%} & 94.2\% & 100.0\% & 0.566 \\
Code-switched & 100.0\% & \textbf{53.6\%} & 98.2\% & 100.0\% & 0.480 \\
\bottomrule
\end{tabularx}
\end{table}

Amount accuracy is essentially perfect in English (99.2\%) and collapses in
Arabic (37.3\%), with code-switched utterances in between (53.6\%, consistent
with many of them expressing the amount in Arabic words even when other parts
of the sentence are English). Type, currency, and date show no comparable
language gap. Section~5.7.2 traces this precisely to how Qwen2.5-7B-Instruct
parses Arabic \emph{spelled-out compound numbers}.

\subsection{End-to-End Audio Pilot (Condition (b))}

The 6-clip pilot (Table~\ref{tab:extract-main}, column (b)) shows perfect
type/amount/date accuracy but weak currency (71.4\%) and category/name
F\textsubscript{1} (0.222/0.333). With $n=6$ (7 gold transactions), each
utterance is worth roughly 14 percentage points, so no strong claim should be
drawn from the absolute numbers---but two of the errors are individually
interesting. One clip's speaker said ``\emph{30 pounds}'' for the price of a
pair of earbuds; British pounds (GBP) are not among MiPay's nineteen supported
currencies (Table~4.4, chosen to cover the Arab region plus major reserve
currencies), so no correct prediction was possible by construction---the gold
label is \texttt{null} and is scored as a currency miss regardless of what the
model outputs. This is a real, useful finding about the currency enum's
coverage, not a model failure (Section~5.7.5). The other systematic pattern in
this pilot is the same category-defaulting behavior seen at full scale
(Section~5.7.3): both purchases in the multi-transaction clip (an iPhone and a
pair of earbuds) were labeled \texttt{shopping} in the gold set but
\texttt{other} by the model.

\subsection{End-to-End Latency}

Table~\ref{tab:latency} reports per-stage and total latency measured directly
from \texttt{POST /transactions/voice} on the 6 recorded clips (5.8--16.8\,s
each; the CPU-only reference machine, no GPU). Figures are the median across
the 6 calls; the longest clip (16.8\,s, two transactions) is an outlier shown
separately.

\begin{table}[H]
\centering
\caption{End-to-end latency on the CPU-only reference machine, measured on the 6-clip audio pilot.}
\label{tab:latency}
\small
\renewcommand{\arraystretch}{1.3}
\begin{tabularx}{\textwidth}{X c c}
\toprule
\textbf{Stage} & \textbf{Median (ms), $n$=6} & \textbf{Longest clip (16.8\,s, 2-tx) (ms)} \\
\midrule
Transcription (Whisper \texttt{small}) & 3{,}890 & 5{,}693 \\
Extraction (Qwen2.5-7B)                & 6{,}995 & 12{,}011 \\
\midrule
\textbf{Total}                         & 11{,}106 & 17{,}936 \\
\bottomrule
\end{tabularx}
\end{table}

Against the NFR-01 budget of $\leq$\,10\,s for a 10-second clip, the median
total (11.1\,s, over clips of 5.8--7.5\,s) modestly \emph{exceeds} the target
on this CPU-only hardware, with extraction---not transcription---as the larger
of the two stages. The 16.8-second, two-transaction clip took 17.9\,s, showing
that longer/compound utterances cost proportionally more at the extraction
stage (larger output, more tokens to decode under the schema constraint) as
well as at the transcription stage. This is consistent with the specification's
own risk assessment (Section~11 of the spec) and is the direct, measured
argument for the GPU-acceleration and smaller-model fallback paths documented
in Section~9.3 of the spec, and for the faster-whisper upgrade path in Chapter~6.

\section{Ablation Studies (Design and Reference Point)}

The three ablations below isolate the contribution of each major configuration
choice. Each requires restarting the containers with a different model and
re-running the full 274-utterance corpus (roughly 30 CPU-minutes per
configuration beyond what was already run for the reference configuration);
running all three was outside the time available for this evaluation pass.
Rather than present tables with blank cells, this section states the exact
comparison each ablation is designed to make, together with the one data
point already measured in each case (the configuration actually deployed and
reported throughout this chapter), so that running the remaining
configurations is a direct drop-in once time allows.

\subsection{Ablation 1 --- Whisper Model Size}

This ablation compares Whisper \texttt{base}, \texttt{small}, and
\texttt{medium} on Arabic, English, and code-switched WER, and on
transcription latency, holding the extraction model fixed. The deployed
configuration, \texttt{small}, is the one measured throughout this chapter:
76.0\% WER on the Arabic pilot ($n=4$), 17.4\% on the English pilot ($n=2$),
and a median transcription latency of 3{,}890\,ms (Table~\ref{tab:latency}).
The expected direction, consistent with Whisper's published scaling
behavior~\cite{radford2023whisper}, is monotonically lower WER moving from
\texttt{base} to \texttt{medium} at the cost of higher latency; \texttt{base}
is the CPU-latency floor and \texttt{medium} the accuracy ceiling for this
pipeline, with \texttt{small} chosen as the balance point for the reference
hardware (Section~1.5.1).

\subsection{Ablation 2 --- Language-Model Size}

This ablation compares \texttt{qwen2.5:3b-instruct} against the deployed
\texttt{qwen2.5:7b-instruct} on amount accuracy, type accuracy, category
F\textsubscript{1}, and extraction latency, holding the speech model fixed.
The deployed 7B configuration measures 65.6\% amount accuracy, 99.3\% type
accuracy, 0.554 category F\textsubscript{1}, and a median extraction latency
of 6{,}995\,ms (Tables~\ref{tab:extract-main} and~\ref{tab:latency}). The
3B model is expected to trade some accuracy--- particularly on category
disambiguation and the compound-number cases discussed in
Section~5.7.2---for lower latency and memory, which is precisely why it is
documented as the low-resource fallback (Section~4.2 of the spec) rather than
the default.

\subsection{Ablation 3 --- Few-Shot versus Zero-Shot Prompting}

This ablation compares the deployed eight-example few-shot prompt against a
zero-shot variant retaining only the system instructions, on amount accuracy,
category F\textsubscript{1}, and transaction-count accuracy. The deployed
few-shot configuration measures 65.6\% amount accuracy, 0.554 category
F\textsubscript{1}, and 100.0\% transaction-count accuracy (i.e.\ every
multi-transaction utterance in the corpus was split into the correct number
of records). The few-shot examples are expected to matter most for exactly
this behavior---multi-transaction splitting and currency disambiguation---since
those are the patterns the examples explicitly demonstrate; a zero-shot
prompt retains the same schema-constrained guarantee of valid output but is
expected to fall back to less consistent splitting on compound utterances.

\section{Qualitative Examples}

Table~\ref{tab:qualitative} shows representative utterances and the pipeline's
\emph{actual} predictions from the evaluation run (not idealized examples),
illustrating the two systematic error modes analyzed in Section~5.7 alongside
cases that work correctly.

\begin{table}[H]
\centering
\caption{Representative utterances and actual observed pipeline output (Arabic transliterated).}
\label{tab:qualitative}
\small
\renewcommand{\arraystretch}{1.3}
\begin{tabularx}{\textwidth}{X X X}
\toprule
\textbf{Utterance (transliterated)} & \textbf{Gold} & \textbf{Predicted} \\
\midrule
``ishtarayt bidala min Carrefour bi-mitayn wa-khamsin gineh al-naharda'' & 250 EGP, groceries, Carrefour & \textbf{150 EGP} (``mitayn'' 200 misread as ``miya'' 100), groceries, Carrefour \\
``da\!fact talatmiya wa-cishrin gineh fi Spinneys'' & 320 EGP, groceries & \textbf{20 EGP} (hundreds word dropped) \\
``ishtarayt ti-shirt min Zara bi-tamanmiya gineh'' & 800 EGP, \textbf{shopping}, Zara & 50 EGP, \textbf{other}, Zara \\
``I bought an iPhone with 100 bucks and earbuds for 30 pounds'' & 100 USD + 30 (currency n/a), \textbf{shopping} & 100 USD + 30, \textbf{other} (both items) \\
``So today I bought \$50 pizza.'' & 50 USD, restaurants, pizza & 50 USD, restaurants, pizza (exact match) \\
``i bought a coffee 50 bucks and a croissant 100 bucks and dad gave me 200'' & 3 transactions: 50/100/200 USD & 3 transactions: 50/100/200 USD (exact match, correct split) \\
``al-jaw hilw al-yawm wal-hamdu lillah'' (no transaction) & \texttt{[]} (empty) & \texttt{[]} (empty, correctly detected) \\
\bottomrule
\end{tabularx}
\end{table}

\section{Discussion and Analysis}

\subsection{Does the System Meet Its Targets?}

Against NFR-02, the system \textbf{exceeds} the type target (99.3\% vs.\
$\geq$\,85\%) and has no explicit target for currency or date but scores 96.6\%
and 99.3\% respectively---all excellent. It \textbf{misses} the amount target
substantially (65.6\% vs.\ $\geq$\,90\%) and the category target
(F\textsubscript{1} 0.554 vs.\ $\geq$\,0.80). Both misses are traced to
specific, fixable causes below rather than being diffuse quality problems,
which changes how they should be read: this is not "the model is generally
unreliable at 65\%,'' it is "the model is unreliable at one specific,
well-defined sub-task (Arabic compound-number parsing) and defaults to
\texttt{other} for six specific, well-defined categories.''

\subsection{Root Cause 1: Arabic Compound-Number Parsing}

The per-language breakdown (Table~\ref{tab:extract-lang}) isolates the amount
failure to Arabic: 37.3\% accuracy versus 99.2\% in English. Manual inspection
of the mismatches (Table~\ref{tab:qualitative}, rows 1--2) shows a consistent
pattern: Qwen2.5-7B-Instruct frequently mis-parses Egyptian-colloquial
\emph{compound} number words---hundreds spoken with a conjunction before tens or
units, e.g.\ ``\emph{talatmiya wa-cishrin}'' (300 + 20 = 320) or ``\emph{mitayn
wa-khamsin}'' (200 + 50 = 250)---dropping the hundreds component entirely
(320~$\rightarrow$~20) or confusing the dual form ``\emph{mitayn}'' (200) with
the singular ``\emph{miya}'' (100), yielding 250~$\rightarrow$~150. Simple
amounts (a single number word, or Western/Arabic-Indic digits) are unaffected;
English and digit-based amounts in any language score near-perfectly. This
error also explains part of the low \texttt{name} score (Section~5.7.2): when
the compound number is mis-parsed, the unrecognized fragment sometimes remains
attached to the \texttt{name} field instead of being cleanly separated (e.g.\
gold \emph{name}=``Awlad Ragab'' vs.\ predicted ``Awlad Ragab
bitalatmiya-wa-tamanya''). This is precisely the class of error the project's
own risk register anticipated (Section~11 of the specification, ``Qwen
mislabels category/currency/amount on rare phrasings''), and it is the strongest
data-backed argument in this thesis for the Chapter~6 upgrade path: either a
larger/newer LLM, or---more directly---extending the deterministic
\texttt{postprocess.py} numeral-normalization layer (already handling
Arabic-Indic digit substitution) with an explicit Arabic
number-word-to-digit converter, so amount extraction no longer depends on the
LLM getting compound number words right.

\subsection{Root Cause 2: A Prompt-Coverage Gap for Six Categories}

The category confusion matrix (computed from the raw prediction capture) shows
that errors are not spread randomly across the seventeen-category taxonomy
in place at evaluation time:
they concentrate overwhelmingly on \texttt{shopping} (19 misclassifications),
\texttt{business} (12), \texttt{gifts\_donations} (10), \texttt{education} (9),
\texttt{entertainment} (9), and \texttt{personal\_care} (8), in every case
predicted as \texttt{other}. Cross-referencing the system prompt
(Appendix~A.4), this is exactly the set of categories for which the prompt
gives \emph{no} worked keyword examples: the prompt enumerates explicit mapping
examples for \texttt{groceries}, \texttt{restaurants}, \texttt{fuel},
\texttt{transport}, \texttt{bills}, \texttt{rent}, \texttt{health},
\texttt{salary}, and \texttt{transfer\_in}, but not for the six categories that
dominate the error count. This is a concrete, low-risk, high-confidence fix:
adding one worked example per uncovered category to the few-shot prompt (no
retraining, no schema change) is the most cost-effective single improvement
identified by this evaluation, and is recommended as an immediate follow-up
ahead of any model upgrade.

\subsection{The \texttt{name} Metric Is Partly an Artifact}

Breaking down the 220 gold/predicted \texttt{name} pairs in the gold condition:
179 (81\%) are cases where the gold label is \texttt{null} (the dataset author
judged no name-worthy detail was present) but the model extracted \emph{some}
value anyway (e.g.\ gold=\texttt{null}, predicted=``McDonald's''); only 30
(14\%) are cases where both gold and prediction contain a value but disagree,
and several of those are themselves downstream contamination from the
compound-number failure (Section~5.7.2) rather than independent naming errors.
Read this way, the raw 28.2\%/67.5\% (accuracy/F\textsubscript{1}) numbers
understate how often the model's \texttt{name} output would actually satisfy a
user reviewing the confirmation screen; the strict exact-match metric simply
does not give credit for a plausible name where the curated gold happened to
leave the field empty. This is reported as a metric limitation, not silently
corrected, so the numbers in Table~\ref{tab:extract-main} remain an honest,
reproducible measurement.

\subsection{The Currency Enum Has a Real, Bounded Gap}

The 6-clip pilot's lowest-scoring field, currency (71.4\%), includes one
genuinely unanswerable item: a speaker said ``pounds'' meaning British pounds,
which is not among MiPay's nineteen supported currencies (the schema targets
the Arab region plus major reserve currencies; Table~4.4). No prediction could
have been marked correct for this item under the current schema. This is
useful, bounded evidence---not of model failure, but of a currency-coverage
decision worth revisiting if the target user base is expected to reference
non-regional currencies.

\subsection{The Role of Human Confirmation}

Every extraction error identified above is caught at the confirmation step
before it reaches the database (Section~3.9, two-step save). The measured
accuracies are therefore a proxy for \emph{how often the user must edit a
field}, not a measure of data integrity: a mis-parsed Arabic amount becomes a
one-tap correction, never a silently corrupted record. Given that the
amount-accuracy shortfall is concentrated in Arabic compound numbers---a
narrow, identifiable linguistic construction rather than a pervasive
weakness---the practical user-facing impact is that Arabic speakers using
three-digit-plus compound amounts will edit the amount field more often than
English speakers, a specific, fixable, and honestly reportable limitation
rather than a systemic reliability problem.

\subsection{Threats to Validity}

The 268-row text corpus is authored by the team and, while constructed to span
the full category/currency/dialect space, may not match the exact distribution
of real user speech. The 6-clip audio pilot is far smaller than the 50--100
clips originally targeted (Section~4.1) and every number derived from it should
be read as indicative, not statistically powered---explicitly flagged wherever
it is cited in this chapter. Two of the six audio clips are repeated takes of
the same underlying sentence (recorded because the first attempts were of poor
quality), which is disclosed rather than hidden, but means the pilot's
effective linguistic diversity is five utterances, not six. Scaling the audio
pilot to the originally targeted size, and re-running it after the two fixes
identified in Sections~5.7.2--5.7.3, are the highest-priority items in the
Chapter~6 future-work list.
